% !TEX root = userguide.tex

\chapter{Introduction}
\pgst{tfemch}

\def\chaptertitle{Introduction}
\def\headdate{5th September 2025}

\section{What is \tfem?}

\tfem\ is a toolkit for the finite element method (\fem).
It is written in \mftn. The core of \tfem\ is moderately sized\footnote{At the
time of writing it contains 44320
lines of \ftn\, excluding comments and empty lines.}
and it is intended to be used in education and
research. It is a general \fem-system, i.e.\ it can
\begin{itemize}
   \item build system of equations for 
         Galerkin \fem, discontinuous Galerkin (\dg) \fem\ and the extended finite element
         method (\xfem),
   \item have varying number degrees of freedom per node,
   \item identify multiple degrees of freedom as a single physical degree of freedom,
   \item group elements, so that different element types 
         and/or different material properties can be used in the same problem,
   \item have layers of degrees, with full freedom of choice in
         which nodes a layer is actually present,
         facilitating implementation of \xfem,
   \item define local transformations of nodal degrees of freedom to apply Dirichlet boundary conditions
         in any direction.
\end{itemize}
Furthermore,
\begin{itemize}
   \item many types of elements are available, including higher-order elements,
   \item has extensive functionality for computing derived quantities,
   \item blended meshes can be defined, which gives the possibility to easily mix different orders of interpolation for mixed or coupled variables.
\end{itemize}
However, the core of \tfem\ has limited scope:
\begin{itemize}
   \item The built-in mesh generator generates structured meshes only.
   \item I/O only consists of writing warning and error messages to standard
         output.
   \item No element subroutines, only modules for
         shape functions and Gauss integration are available.
   \item No graphics routines, although some basic post-processing such as
         sampling vectors and integration
         on curves, surfaces and objects, is available.
   \item No linear equation solvers\footnote{Interfaces to external solvers are available 
         as an add-on. For example, 
   there are interfaces to the MA41 (unsymmetric
         matrices) and MA57 (symmetric matrices) sparse direct solvers
         of the HSL library \cite{HSL} available in the add-on hsl, see Sec~\ref{sec:hsl_addon}.}.
   \item No high level routines such as time stepping or
         non-linear equation solvers are available.
\end{itemize}
So the core of
\tfem\ provides basic functionality that every \fem\ program needs, but
additional functionality is necessary to make a complete program. This
functionality can be provided through so-called add-ons along with the core 
routines of \tfem, but this is not necessary.
For example, a PhD student may decide to write his/her own code for that and
link his program to the core \tfem\ modules only. A wide range of 
add-ons already exist\footnote{At the time
 of writing all add-ons together have 63567
lines of \ftn\, excluding comments and empty lines.}
such as elements for (viscoelastic) flow problems, interfaces
to (parallel) external solvers and
readers for external mesh generators\footnote{Several are available, including a reader for meshes made by
Gmsh.}.
At least one add-on for graphics must be available to obtain a functional
program\footnote{For 2D graphics and simple 3D graphics the add-on figplot
is available.
For more advanced graphics add-ons that can write to 
VTK-files, gmsh-files or Tecplot-files, are available.
VTK-files can be read by various freely available
visualization programs, of which Paraview is the most important one.
Most examples in the userguide use either figplot or Paraview. Gmsh postprocessing is used for solid mechanics problems using trusses and beam structures.}.

A \tfem\ \emph{experimental} repository is available that
contains add-ons and examples that are not (yet) fit to include in the regular
\tfem\ repository, but can still be useful.

Add-ons and examples that are difficult to maintain or are replaced by better
ones, are moved to a separate \tfem\ \emph{legacy} repository. In this way, these add-ons and
examples are still available if needed.

\tfem\ has been loosely modeled after the part of \sep\ \cite{Segal93} 
that is run from a main program and calls several, not too high-level,
subroutines.
\tfem\ is controlled from a main \ftn\ program and doesn't need any
input from a file. Like \sep\ it has two basic structures, the mesh and the
problem, but in \tfem\ these structures are stored using modern type
definitions in \ftn. People familiar with \sep, will sometimes feel the spirit
of \sep\ when using \tfem, although the details are quite different.

\section{Getting started}
\label{sec:gettingstarted}

We assume that you have a working compiler 
environment and have a \tfem-tree somewhere on your (Linux/Unix or macOS) system
that has been compiled\footnote{The details of setting up a
compiler environment and compiling \tfem\ are not complicated but beyond the scope of this guide.
See the directory \texttt{install} of the \tfem\ tree for more information.
\tfem\ is not supported on Windows natively, but running Linux in a Docker container or under Windows Subsystem for Linux (WSL) are good alternatives.}.

First you have to add the \texttt{bin}-subdirectory of the \tfem\ root directory
to your path. Test it by typing 
\begin{quote}
  \texttt{which fmm}
\end{quote}
at the command line and see whether the correct path is used to the
\texttt{fmm} command.

Now you can develop your own main program, but at this point it is easier to
copy an example. For that, type
\begin{quote}
   \texttt{getexample poisson}
\end{quote}
which creates the (sub-)directory \texttt{poisson} containing the example files\footnote{The command \texttt{getexample} has an option 
(\texttt{-M})
to generate the \texttt{Mdefs.mk} and \texttt{Makefile} files and an option 
(\texttt{-b}) to 
completely build an example, but it is
advised not to use these options before you have become familiar with all the basic
steps of building an example.}. Now 
\texttt{cd} to that directory.
You will see several source files (extension \texttt{.f90}) in your directory.
Among these are several main programs.

In order to compile all files and link the main program(s), type the following
commands
\begin{quote}
   \texttt{getmdefs >Mdefs.mk}
\end{quote}
which will create the file \texttt{Mdefs.mk}, and
\begin{quote}
   \texttt{fmm >Makefile}
\end{quote}
which creates the file \texttt{Makefile}.
There are some restrictions on the use of
\texttt{fmm}, which are described in Sec.~\ref{sec:fmm}. In particular: main
programs must have an \texttt{end program} statement and modules
must be in separate files having names identical to the module name. The
following command removes the module restriction, but is slower:
\begin{quote}
   \texttt{fmme >Makefile}
\end{quote}
The actual compilation/linking is done by typing the command
\begin{quote}
   \texttt{make}
\end{quote}
The \texttt{make} utility uses the file \texttt{Makefile} to find the
dependencies between the various source files. So if you edit a source file 
and type \texttt{make} again, 
all source files that depend on the changed file
will be recompiled. Also all main programs
will be recompiled/linked that depend on the recompiled files. The file
\texttt{Mdefs.mk} is included by \texttt{Makefile} and
contains the definitions that are used to compile/build \tfem\ and normally
doesn't need to be changed. If you change the dependencies between the source
files, for example by adding a source file or adding an explicit
dependency\footnote{In \ftn\ an explicit dependency is created by 
adding a \texttt{use <module>} statement, where
\texttt{<module>} has been defined in your own directory.}, the
file \texttt{Makefile} needs to be recreated using \texttt{fmm}.
The files produced by \texttt{make} can be removed by
\begin{quote}
   \texttt{make clean}
\end{quote}
and you return to the situation before typing \texttt{make}.
The \texttt{make} utility is available on all Linux/Unix and Unix-like systems.

Many examples in this User Guide use the add-on figplot for plotting (see
Sec.~\ref{sec:figplot}). Figplot produces files in the Fig-format as
used by the program \texttt{xfig} (not to be confused with the fig-files of
Matlab). Several shell scripts have been written to view and/or
convert\footnote{The fig-files can also be viewed/edited using the program
\texttt{xfig}, of course.} the fig-files.

The command 
\begin{quote}
   \texttt{viewfig [figfile(s)]}
\end{quote}
shows on the screen the figfiles given or all figfiles if no files are
given\footnote{The \texttt{viewfig} command converts the fig-files to eps-files
and creates a single pdf file to view them all at once. A disadvantage
is, that it produces full pages and not cropped ones.
The command 
\begin{quote}
   \texttt{viewfigpdf [figfile(s)]}
\end{quote}
is similar to \texttt{viewfig}, but converts the fig-files directly to pdf-files.
In does crop the pages, but is slower.
The command 
\begin{quote}
   \texttt{viewfigpng [figfile(s)]}
\end{quote}
converts the fig-files directly to png-files, which are viewed
using an image viewer.
}.

The command \texttt{fig2} converts fig-files to another graphics format, for example
\begin{quote}
   \texttt{fig2 pdf [figfile(s)]}
\end{quote}
will generate pdf-files of the figfiles given
or all figfiles if no files are given. Typing just:
\begin{quote}
   \texttt{fig2}
\end{quote}
will give an overview of all the graphics formats available.

\section{Modern Fortran}

\tfem\ has been written in the computer language (modern) \ftn\footnote{Modern Fortran usually means
Fortran~90 and later. For \tfem\ we require modern Fortran, but don't restrict to a particular version of
the standard. We follow the latest standard (currently \ftn~2023), provided the particular
feature of the standard is supported by the three most
important compilers (NAG Fortran, Intel Fortran, GNU Fortran).}. Although the first
Fortran program was written more than 65 years (!) ago, nowadays it is a very
modern language, designed for the professional computational scientist:
\begin{itemize}
   \item It is very efficient.
   \item It has an array/matrix notation similar to Matlab.
   \item It has parallel programming syntax built in (coarrays\footnote{
         Coarrays are part of the \ftn~2008 standard, but not yet fully implemented in
         all compilers.}).
   \item It is available for every relevant computer system and compilers are available from various
         vendors\footnote{There is also an excellent open-source compiler
         (\texttt{gfortran}), which is part of gcc. Also Intel Fortran (\texttt{ifx}) has been included in the oneAPI toolkit and is free to use.}.
   \item Old, but usually very well tested, code written in old-style
         Fortran can still be compiled and used in \ftn.
\end{itemize}
In order to use \tfem, the user needs to have
some knowledge of \ftn. How much depends
of course on how you are going to use it. Do you plan on writing new code or
only change an existing main program slightly? In the latter case, you might
even get away with knowing only how to continue a statement on the next line. However, it is recommended to
go to \href{https://fortran-lang.org}{fortran-lang.org}, where users are building a community around the \ftn\ language. Click on \boxed{\texttt{Get Started}} and start reading the Quickstart tutorial.
A good book for learning
\ftn\ is the book\footnote{The book is available as an e-book from the library of the TU/e.} of Brainerd \cite{Brainerd2015}.
For those who are already familiar with \mftn, but need a good reference, the
book of Metcalf, Reid, Cohen \& Bader \cite{metcalf_modern_2024} is highly recommended.

\section{Editing source code}

For editing source code you need an editor that, preferably, can
handle Fortran syntax highlighting and has a feature to quickly traverse through the source code.

Most of \tfem\ has been developed using Vim
(\texttt{http://www.vim.org}), which is available for all systems. 
Vim can use the \texttt{tags} files to quickly move around in the source files (see Sec.~\ref{sec:viewing}). Vim is very powerful, but has a learning curve in order to get used to the commands. 
Type \texttt{vimtutor} at the command line to get a half an hour introduction course. 

Visual Studio Code (\texttt{https://code.visualstudio.com}) is also a powerful editor with a nice GUI interface that has become quite popular in the programming community. It also has good Fortran support with the ``Modern Fortran'' extension.

\section{Viewing source code}
\label{sec:viewing}

A reference guide for \tfem\ does not exist.
Therefore, to know all the
possibilities of various routines, it is important to view the (comments) in
the source. To make that an easy process, we use so-called tags.

When compiling/building \tfem,
but also when using \texttt{make} on your own files,
a file \texttt{tags} is made. Such a file contains hooks to various points in
the source code. Various editors (including Vim)
can read this file and provide a 
mechanism to immediately jump to a particular point in a source file.

\tfem\ has a command \texttt{tfvi} to jump to a tag, using Vim. So
\begin{quote}
   \texttt{tfvi problem\_t}
\end{quote}
will jump to the definition of the derived type \texttt{problem\_t} and 
\begin{quote}
   \texttt{tfvi problem\_definition}
\end{quote}
will jump to the first line of the subroutine \texttt{problem\_definition}.

You
can also use commands within an editor to jump to a tag. So while editing
or viewing with Vim, you can type for example
\begin{quote}
   \texttt{:tag mesh\_t}
\end{quote}
to jump to the tag \texttt{mesh\_t}. It you don't know exactly the tag name:
it is possible to use name expansion, where \texttt{TAB} will expand the tag.
Typing \texttt{Ctrl-D} will give all the remaining possibilities.

Another possibility to jump to a tag is to
put your cursor on a word in a file and type
\texttt{Ctrl-]} to jump to the tag defined by the underlying word\footnote{Sometimes there are multiple possibilities and \texttt{Ctrl-]} will jump to the first possibility by default. To get an overview of all possibilities use \texttt{g-]}, instead.}.
You can type as many times \texttt{Ctrl-]} as you want to jump to new
positions. The previous positions are stored (on a so-called stack) and
you can always go back by typing \texttt{Ctrl-t} to go
backwards in the editing/viewing stack. NOTE: Vim has the annoying feature that 
it also reads the tags file in the directory of the current file leading to
multiple entries for the same tag when the stack is larger than one. To switch
off this feature forever put 
\begin{quote}
   \texttt{set tags=tags}
\end{quote}
or
\begin{quote}
   \texttt{set cpoptions+=d}
\end{quote}
in your \texttt{\$HOME/.exrc} of \texttt{\$HOME/.vimrc} file.

Note that if your \texttt{make} command is successful, a \texttt{tags} file is
made in your own directory and you can jump in the same way through your own
source. If \texttt{make} is not successful you can force it to make a
\texttt{tags} file by
\begin{quote}
   make tags
\end{quote}

\section{Documentation}

Type
\begin{quote}
   \texttt{tfdoc}
\end{quote}
to see a documentation list. A document can be viewed with
\begin{quote}
   \texttt{tfdoc <document>}
\end{quote}
An alias for \texttt{tfdoc userguide} is
\begin{quote}
   \texttt{tfug}
\end{quote}
and will show this userguide on your screen.

Element modules are available as add-ons and the required coefficients are
documented using \texttt{README\_coefficients\%i} and 
                 \texttt{README\_coefficients\%r} text files in the add-on
directory. 
All the available add-ons having these files can be listed by
\begin{quote}
   \texttt{showcoef -l}
\end{quote}
Showing these files on the screen can be achieved by
\begin{quote}
   \texttt{showcoef [add-on]}
\end{quote}
or 
\begin{quote}
   \texttt{showcoef [add-on] | less}
\end{quote}
to easily scroll through the file.
Omitting the add-on gives the coefficients of all add-ons.

\newpage
\section{Getting \tfem}

The source code of \tfem, including the documentation, is stored in a Git
repository on the GitLab server \href{https://gitlab.tue.nl}{\texttt{https://gitlab.tue.nl}}. You can login directly to the GitLab server to view source files, browse the commit history, report an issue, create a merge request and more. However, if you want to use \tfem, you need to obtain a local copy and compile the source.

You can copy (`clone') a fresh \tfem\ tree from the GitLab server by typing at the command line
\begin{quote}
   \texttt{git clone https://gitlab.tue.nl/mahulsen/tfem.git tfem}
\end{quote}
This command will create a subdirectory \texttt{tfem} that contains the \tfem\ tree. See the file
\texttt{README.md}
for further directions on building \tfem\ from the source.

To see whether \tfem\ has been updated, cd to the \tfem\ tree and type
\begin{quote}
   \texttt{git remote show origin}
\end{quote}

You can update to the latest revision of \tfem\ by going to the root of the
\tfem\ tree and just type
\begin{quote}
  \texttt{git pull} 
\end{quote}
It is advised to always
rebuild the whole tree after an update to get rid of possible
incompatibilities between new
and previously compiled routines.

To see a log of commits type
\begin{quote}
   \texttt{git log}
\end{quote}
which is automatically piped through \texttt{less}\footnote{I had to set the environment variable \texttt{export LESS=R} to have a proper output with colors.}.

If you find the command line interface of Git difficult to use, maybe \href{https://desktop.github.com}{GitHub Desktop} is something for you.

Further information on the use of Git is available at \href{https://git-scm.com}{\texttt{https://git-scm.com}}.
